%=============================================================
%  Module I.4 / L10 — Scattering from a Potential Barrier
%  Series I > Module I.4 > Lecture L10
%  Duration: 90 minutes  |  All tiers (T1–T5)
%=============================================================
\renewcommand{\lectureCode}{I.4 / L10}

\section*{L10 — Scattering from a Potential Barrier}
\label{sec:L10}
\addcontentsline{toc}{section}{L10 — Scattering from a Potential Barrier}

%--- Lecture header
\begin{tcolorbox}[lecturebox, title={Module I.4 / L10 — Metadata}]
\begin{tabular}{@{}ll@{}}
\textbf{Series:}    & I — Introductory Quantum Mechanics \\
\textbf{Module:}    & I.4 — Paradigm Physical Systems \\
\textbf{Lecture:}   & L10 — Scattering from a Potential Barrier \\
\textbf{Duration:}  & 90 minutes \\
\textbf{Tiers:}     & T1 (HS) · T2 (BegUG) · T3 (AdvUG) · T4 (MSc) · T5 (PhD) \\
\textbf{Preceding:} & I.4 / L09 — The Hydrogen Atom \\
\textbf{Following:} & I.4 / L11 — Scattering from a Potential Well \\
\textbf{Prerequisites:} & Modules I.1–I.3 complete \\
\end{tabular}

\medskip
\textbf{Key concepts:} Rectangular barrier, transfer matrix, tunnelling, Gamow factor, above-barrier reflection
\end{tcolorbox}

%------------------------------------------------------------
\subsection*{L10.0 — Learning Objectives (per tier)}
\label{sec:L10.0}

\tiertag{tier1col}{Tier 1 — HS}\quad
Identify the main physical phenomenon studied in this lecture; describe it in
non-mathematical terms; state one real-world application.

\tiertag{tier2col}{Tier 2 — BegUG}\quad
Apply the key equations to compute numerical results; verify consistency with
known limits; translate between physical and mathematical descriptions.

\tiertag{tier3col}{Tier 3 — AdvUG}\quad
Derive central results from first principles; prove key identities; identify
the mathematical structure underlying the physical phenomenon.

\tiertag{tier4col}{Tier 4 — MSc}\quad
Analyse formal extensions; connect to symmetry principles; generalise to
related systems; identify the role of this lecture in the broader programme.

\tiertag{tier5col}{Tier 5 — PhD}\quad
Establish rigorous mathematical foundations; identify open research problems;
connect to modern literature; critique standard textbook treatments.

%--- Lecture content -----------------------------------------

\subsection*{L10.1 — Setup and Transfer Matrix}
Energy $E$ incident on rectangular barrier $V_0$, width $a$.
Three regions; wavevectors $k_1=\sqrt{2mE}/\hbar$, $k_2=\sqrt{2m|E-V_0|}/\hbar$.

\begin{tcolorbox}[lecturebox, title={Transmission and Reflection}]
Transmission: $T = |t|^2$; reflection: $R = |r|^2$; $T+R=1$ (probability current).
Resonance (above barrier): $k_2 a = n\pi$, $T=1$.
Tunnelling ($E<V_0$): $\kappa = \sqrt{2m(V_0-E)}/\hbar$,
\begin{equation}
  T = \left[1+\frac{(k_1^2+\kappa^2)^2}{4k_1^2\kappa^2}\sinh^2(\kappa a)\right]^{-1}
  \approx e^{-2\kappa a}\text{ (thin barrier)}.
  \label{eq:L10.1}
\end{equation}
\end{tcolorbox}

\subsection*{L10.2 — Gamow Factor and Applications}
$G = -2\int_{x_1}^{x_2}\kappa(x)\,dx$ (WKB); $T\approx e^G$.
Alpha decay: $G = \pi Z_1 Z_2 e^2/\hbar v$ (Coulomb barrier).
Applications: STM, alpha decay, nuclear fusion, Josephson junction.


%------------------------------------------------------------
\subsection*{L10.C — Connections}
\label{sec:L10.C}
\textbf{Backward:} I.4/L09 — The Hydrogen Atom and Modules I.1–I.3.
\textbf{Forward:} I.4/L11 — Scattering from a Potential Well builds directly on the formalism developed here.
\textbf{Cross-module:} Series~II modules extend these results to many-body and
advanced settings; Series~III applies them to molecular electronic structure.

%=============================================================
%  ASSESSMENT BUNDLE — L10
%=============================================================
\newpage
\section*{L10 — Assessment Artifact Bundle}

\begin{tcolorbox}[ugconceptbox, title={SCQU · L10 · Tiers 1–3 · 15–20 min}]
\textit{Ten simple concept questions (mix MC/TF/short) targeting the core results
of Scattering from a Potential Barrier. See artifact file \texttt{L10_artifacts.tex} for full questions.}

\textbf{Rubric:} 2 pts each (1.5 correct + 0.5 justification). Total: 20 pts.
\end{tcolorbox}

\begin{tcolorbox}[ugconceptbox, title={CCQU · L10 · Tiers 2–3 · 25–35 min}]
\textit{Ten complex concept questions: ``explain why,'' ``what would change if,''
``identify the flaw.'' Full questions in \texttt{L10_artifacts.tex}.}

\textbf{Rubric:} 2 pts each (1 key point + 1 argument). Total: 20 pts.
\end{tcolorbox}

\begin{tcolorbox}[ugprobbox, title={SPSU · L10 · Tiers 1–3 · 60–90 min}]
\textit{Ten problems (Part A: mechanics, Part B: interpretation, Part C: translation).
Full problems in \texttt{L10_artifacts.tex}.}

\textbf{Rubric:} Result 60\%, Method 30\%, Interpretation 10\%.
\end{tcolorbox}

\begin{tcolorbox}[ugprobbox, title={CPSU · L10 · Tier 3 · 3–6 hrs (4-layer)}]
\textit{Five multi-part derivation problems with four-layer scaffolding.
Layers 1–4 in \texttt{L10_ex01\_problem.tex}, \texttt{hints}, \texttt{adv\_hints}, \texttt{solution}.}

\textbf{Rubric:} Assumptions 15\%, Proof structure 45\%, Algebra 25\%, Insight 15\%.
\end{tcolorbox}

\begin{tcolorbox}[ugprobbox, title={SPJU + CPJU · L10 · Tiers 2–3}]
\textit{Simple project (2–6 hrs): structured computational/analytical task.
Complex project (8–15 hrs): open-ended synthesis.
Full specs in \texttt{L10_artifacts.tex}.}
\end{tcolorbox}

\begin{tcolorbox}[gradconceptbox, title={SCQG + CCQG · L10 · Tiers 4–5}]
\textit{Ten simple + ten complex concept questions at graduate level (rigour, domain
conditions, named theorems). Full questions in \texttt{L10_artifacts.tex}.}
\end{tcolorbox}

\begin{tcolorbox}[gradprobbox, title={SPSG + CPSG · L10 · Tiers 4–5}]
\textit{Ten simple + five complex graduate problems.
CPSG with four-layer format. See \texttt{L10_artifacts.tex}.}
\end{tcolorbox}

\begin{tcolorbox}[gradprobbox, title={SPJG (×5) + CPJG (×5) · L10 · Tiers 4–5}]
\textit{Five simple (3–8 hrs each) and five complex (10–20 hrs each) graduate projects.
Full specs in \texttt{L10_artifacts.tex}.}
\end{tcolorbox}

\begin{tcolorbox}[researchbox, title={WRQ (×5) + ORQ (×5) · L10 · Tiers 4–5}]
\textit{Five well-defined research questions (5–10 hrs each) and five open-ended
research questions (10–20 hrs each). Full prompts in \texttt{L10_artifacts.tex}.}
\end{tcolorbox}
